%% Generado por gen.py desde consolidado.json y las Bases.
La Tabla~\ref{tab:anx-vacios} presenta vacíos e inconsistencias para revisión.\par
\begin{tablalafrox}{Vacíos e inconsistencias para revisión. Fuente: contraste de Bases, SD2 y SD3.}{tab:anx-vacios}%
  {L{1.5cm} Y L{10cm}}%
  {\cab{ID} & \cab{Materia} & \cab{Nota de trabajo: evidencia requerida}}
  V-01 & Número de instalaciones: cinco en entrevista, seis en volumetría; SD2 indica 14 CD. & Confirmar universo de instalaciones; no trasladar 14 CD al diseño. \\
  V-02 & La ponderación del Formulario T-21 suma 98 \% aunque su total declara 100 \%. & Contrastar el planteamiento original y documentar tratamiento antes de cerrar la consulta. \\
  V-03 & 84 conductores y peonetas frente a 42 conductores propios en la Tabla 14.1. & Contrastar el planteamiento original y documentar tratamiento antes de cerrar la consulta. \\
  V-04 & Cuatro ambientes (Art. 24 y E-25) frente a cinco más recuperación (Art. 3 y RT-04.01). & Contrastar el planteamiento original y documentar tratamiento antes de cerrar la consulta. \\
  V-05 & Códigos de sincronización y otras materias del capítulo 15 no coinciden con BTT. & Usar código transversal por materia y registrar el valor particular del caso; completar cotejo por fila. \\
  V-06 & Relato de cortes de 2 h frente a la exigencia de 14 h. & Contrastar el planteamiento original y documentar tratamiento antes de cerrar la consulta. \\
  V-07 & Alta y recuperación de credenciales de personas sin correo. & Contrastar el planteamiento original y documentar tratamiento antes de cerrar la consulta. \\
  V-08 & 41 \% de recepciones del producto retirado sin lote en origen. & Contrastar el planteamiento original y documentar tratamiento antes de cerrar la consulta. \\
  V-09 & No existe saldo inicial de envases por cliente. & Contrastar el planteamiento original y documentar tratamiento antes de cerrar la consulta. \\
  V-10 & El Cap. 5 manda reemplazar la planilla de reposición y el catálogo inicial no tenía requerimiento. & Contrastar el planteamiento original y documentar tratamiento antes de cerrar la consulta. \\
  V-11 & Faltaban la hora real de entrega y el monto esperado de rendición, sin los cuales el OTIF y la cuadratura no se calculan. & Contrastar el planteamiento original y documentar tratamiento antes de cerrar la consulta. \\
  V-12 & Regla escrita de excursión térmica y responsable de la disposición. & Contrastar el planteamiento original y documentar tratamiento antes de cerrar la consulta. \\
  V-13 & Inicio calendario del contrato no consolidado. & Confirmar fecha de inicio y compatibilidad de los meses 16 y 21 con congelamientos y enero de 2029. \\
  V-14 & La suspensión del proveedor de lácteos vence en septiembre de 2026, antes del inicio del contrato; no está definido qué evidencia de trazabilidad acepta el proveedor. & Contrastar el planteamiento original y documentar tratamiento antes de cerrar la consulta. \\
  V-15 & Disponibilidad y retiro del planificador. & Confirmar fecha; conservar talleres meses 1 a 3 del SD2, sin fijar retiro entre meses 8 y 13. \\
  V-16 & Condiciones contractuales de acceso a los datos de la telemetría del tercero. & Confirmar acceso, condiciones contractuales e interfaces antes de definir la integración. \\
  V-17 & Rotación de 38 \% atribuida a conductores externos en SD2. & El caso atribuye ese porcentaje a preparación de pedidos; no trasladarlo a conductores. \\
  V-18 & SD2 menciona 99,5 \% y referencias RNF que requieren cotejo. & Definir disponibilidad por servicio contra Bases; no usar experiencia corporativa como SLA contractual. \\
  V-19 & SD2 propone ventanas comerciales y corte horario. & Confirmar capacidad y condiciones de la promesa antes de convertirla en garantía de servicio. \\
  V-20 & Referencias a arquitectura, EDT, pruebas e innovaciones de documentos aún no revisados. & Completar vínculos tras revisar esos documentos; mantener campos vacíos. \\
  V-21 & SD2 menciona SAP Business One y Manhattan; no se adopta la marca como evidencia técnica. & Acreditar fabricante, versión e interfaces antes de diseñar integración o migración. \\
\end{tablalafrox}
